\documentclass[11pt,a4paper]{article}
\usepackage[T1]{fontenc}
\usepackage[utf8]{inputenc}
\usepackage{amsmath,amssymb,amsthm}
\usepackage[margin=27mm]{geometry}
\usepackage[colorlinks=true,urlcolor=blue,citecolor=blue]{hyperref}
\newtheorem{theorem}{Proposition}
\newtheorem{corollary}[theorem]{Corollaire}
\newcommand{\tr}{\operatorname{Tr}}
\newcommand{\ket}[1]{\lvert#1\rangle}
\newcommand{\bra}[1]{\langle#1\rvert}
\newcommand{\norm}[1]{\left\lVert#1\right\rVert}
\title{Piste 51 : coût exact de préparation d'ensembles\\avec accès partiel à un environnement}
\author{Note de recherche OQM --- dérivations assistées par IA}
\date{17 novembre 2026}
\begin{document}
\maketitle
\begin{abstract}
Nous calculons un défaut de réalisation de toutes les décompositions d'une
marginale de qubit dans un modèle pur à deux branches. Le résultat est
$s_B=(1-\kappa)\sqrt{pq}$, où $\kappa$ est le recouvrement des branches
inaccessibles. Une cible binaire et une mesure explicite atteignent le minimax.
Dans une chaîne à rotations contrôlées sommables, cela donne des bornes
supérieure et inférieure concordantes et un coût spatial exponentiel calculable.
Il s'agit d'une application de modèles et d'outils connus : la priorité de cette
formule n'est pas établie. Aucun problème ouvert sur une phase topologique
ni aucune réciproque quantitative générale de dualité de Haag n'est résolu.
\end{abstract}

\section{Tâche et conventions}
Pour un état $\rho_{AB}$, fixons $\rho_A$. Un ensemble cible est une famille
finie $\lambda_i\geq0$ de somme $\rho_A$, avec probabilités incluses.
Nous étudions
\begin{equation}\label{def}
s_B(\rho)=\sup_{\{\lambda_i\}:\,\sum_i\lambda_i=\rho_A}
\inf_{\{E_i\}\text{ POVM sur }B}
\frac12\sum_i\norm{\lambda_i-\tr_B[(I\otimes E_i)\rho_{AB}]}_1.
\end{equation}
Les cibles et mesures ont les mêmes étiquettes. Les mesures peuvent dépendre
de la cible. On ne mesure ni la robustesse de steering contre les modèles à
états cachés, ni une tâche à plusieurs réglages simultanés. Le ``complete
steering'' de Jevtic et al.\ \cite{ellipsoid} porte sur les décompositions
\emph{à l'intérieur de l'ellipsoïde accessible}; ici le supremum inclut
toutes les décompositions de $\rho_A$, y compris hors de cet ellipsoïde.

\section{Un minimax exactement résolu}
Considérons l'état pur
\begin{equation}\label{state}
\ket\Psi=\sqrt p\ket0_A\ket{b_0}_B\ket u_C
+\sqrt q\ket1_A\ket{b_1}_B\ket v_C,
\qquad p+q=1,\quad 0<q\leq p,
\end{equation}
où $b_0,b_1$ sont orthonormaux et $u,v$ unitaires. Le système $C$ est
inaccessible. Absorbons une phase dans $b_1$ de sorte que
$\langle v|u\rangle=\kappa\in[0,1]$. Sur le support logique de Bob,
\[
\rho_{AB}=p\ket{00}\bra{00}+q\ket{11}\bra{11}
+\kappa\sqrt{pq}(\ket{00}\bra{11}+\ket{11}\bra{00}).
\]
La compression d'une POVM physique sur ce support est une POVM de qubit ;
réciproquement toute POVM logique s'étend en attribuant le complément
orthogonal à une issue. Il n'y a donc pas de restriction cachée des mesures.

\begin{theorem}\label{exact}
Pour l'état \eqref{state},
\[
\boxed{s_B(\rho)=(1-\kappa)\sqrt{pq}.}
\]
Le supremum est atteint par les deux opérateurs
$\lambda_\pm=\ket{\phi_\pm}\bra{\phi_\pm}/2$,
où $\ket{\phi_\pm}=\sqrt p\ket0\pm\sqrt q\ket1$.
\end{theorem}
\begin{proof}
Écrivons $D=\operatorname{diag}(\sqrt p,\sqrt q)$ et
$\lambda_i=\begin{pmatrix}x_i&z_i\\\bar z_i&y_i\end{pmatrix}$.
La POVM logique $E_i=(D^{-1}\lambda_iD^{-1})^T$ est positive et de somme $I$.
Elle prépare
\[
\mu_i=\begin{pmatrix}x_i&\kappa z_i\\\kappa\bar z_i&y_i\end{pmatrix}.
\]
Ainsi son coût est $(1-\kappa)\sum_i|z_i|$. La positivité puis
Cauchy--Schwarz donnent
$\sum_i|z_i|\leq\sum_i\sqrt{x_i y_i}\leq\sqrt{pq}$.
Cette construction fonctionne pour tout nombre fini d'issues.

Pour la borne inverse, prenons la cible binaire annoncée. Tout effet
$0\leq E\leq I$ satisfait $|E_{01}|\leq1/2$, car $\norm{2E-I}\leq1$.
La différence $\Delta=\lambda_+-\mu_+$ a donc
\[
\norm\Delta_1\geq\tr(X\Delta)
=\sqrt{pq}-2\kappa\sqrt{pq}\,\Re E_{10}
\geq(1-\kappa)\sqrt{pq},
\quad X=\begin{pmatrix}0&1\\1&0\end{pmatrix}.
\]
Les deux différences sont opposées puisque les sommes valent $\rho_A$.
Le coût total est donc $\norm\Delta_1$. Les effets
$E_\pm=(I\pm X)/2$ atteignent l'égalité.
\end{proof}

\section{Comparaison au défaut d'espérance conditionnelle}
Définissons $\delta_*=\inf_\tau\frac12\norm{\rho_{AC}-\rho_A\otimes\tau}_1$.
C'est le défaut optimal de préservation de l'état des espérances
$E_\tau(X\otimes Y)=X\tr(\tau Y)\otimes I$ dans le cadre matriciel.
Ici
\[
\rho_{AC}=p\ket0\bra0\otimes\ket u\bra u
+q\ket1\bra1\otimes\ket v\bra v.
\]
En notant $U=\ket u\bra u$, $V=\ket v\bra v$, on obtient
\begin{align*}
2\delta_\tau&=p\norm{U-\tau}_1+q\norm{V-\tau}_1\\
&\geq q\norm{U-V}_1+(p-q)\norm{U-\tau}_1.
\end{align*}
L'égalité pour l'infimum est atteinte en $\tau=U$, d'où
\begin{equation}\label{delta}
\boxed{\delta_*=q\sqrt{1-\kappa^2},\qquad
s_B=\sqrt{pq}\left(1-\sqrt{1-(\delta_*/q)^2}\right).}
\end{equation}
Si la fidélité de l'espérance est exigée, approcher $U$ par des états
fidèles conserve cet infimum. La seconde égalité suppose $q$ fixé et
$0\leq\delta_*\leq q$.

Pour $\kappa=0$ et $q\to0$, on retrouve
$s_B=\sqrt{\delta_*(1-\delta_*)}$ : aucune borne linéaire uniforme
en $\delta_*$ n'est possible. À $q>0$ fixé et $\kappa\to1$, au contraire,
\[
s_B\sim\frac{\sqrt{pq}}{2q^2}\delta_*^2.
\]
Ces deux limites ne commutent pas dans une estimation uniforme : la borne
générale en racine carrée peut être très pessimiste en régime de faible fuite.

\section{Accès spatial croissant : taux avec borne inférieure}
Plaçons un qubit $A$, un qubit drapeau $B_0$ et des qubits $E_j$, $j\geq1$.
Partons de $(\sqrt p\ket{00}+\sqrt q\ket{11})_{AB_0}$ et de
$\bigotimes_j\ket0_{E_j}$. Sur $A E_j$, appliquons
\[
U_j=\ket0\bra0_A\otimes I+\ket1\bra1_A\otimes R(\theta_j),
\quad R(\theta)=\begin{pmatrix}\cos\theta&-\sin\theta\\
\sin\theta&\cos\theta\end{pmatrix}.
\]
Supposons $0\leq\theta_j<\pi/2$ et $\sum_j\theta_j<\infty$.
Les $U_j$ commutent et $U=\prod_jU_j$ converge en norme dans l'algèbre
quasi-locale, puisque $\norm{U_j-I}=2\sin(\theta_j/2)\leq\theta_j$.
L'état obtenu est pur. Bob contrôle $B_R=B_0E_1\cdots E_R$ ; le reste
forme $C_R$. Le drapeau rend les deux branches accessibles orthogonales,
même si les branches de $E_1\cdots E_R$ ne le sont pas.

\begin{corollary}\label{spatial}
Avec $V_R=\sum_{j>R}\theta_j^2$,
\begin{align}
s_{B_R}&=\sqrt{pq}\left(1-\prod_{j>R}\cos\theta_j\right),\label{spatialexact}\\
\sqrt{pq}(1-e^{-V_R/2})&\leq s_{B_R}\leq\frac{\sqrt{pq}}2 V_R.\label{bounds}
\end{align}
En particulier $s_{B_R}\sim\sqrt{pq}V_R/2$ si $V_R\to0$ avec
$V_R>0$. Si la queue est nulle, le défaut est exactement nul.
\end{corollary}
\begin{proof}
Le recouvrement des branches inaccessibles vaut le produit dans
\eqref{spatialexact}. Appliquer la proposition \ref{exact}.
Pour une suite finie, $1-\prod a_j\leq\sum(1-a_j)$ pour
$0\leq a_j\leq1$. Utiliser $1-\cos\theta\leq\theta^2/2$ et
$\cos\theta\leq e^{-\theta^2/2}$ sur $[0,\pi/2)$, puis passer à la
limite des produits. Le quotient des deux bornes tend vers un.
\end{proof}

Pour $\theta_j=\theta_0r^j$, $0<r<1$, $0<\theta_0<\pi/2$,
\begin{equation}\label{rate}
s_{B_R}\sim\frac{\sqrt{pq}\theta_0^2}{2(1-r^2)}r^{2(R+1)}.
\end{equation}
La troncature d'unitaire, elle, vérifie
$\norm{U-U_{\leq R}}\leq\theta_0r^{R+1}/(1-r)$.
La norme d'opérateur contrôle une erreur au premier ordre ; la tâche
d'ensembles présente ici une décroissance du second ordre.
Pour garantir $s_{B_R}\leq\varepsilon$, il suffit de choisir
\[
R\geq\max\left\{0,\left\lceil
\frac{\log[\sqrt{pq}\theta_0^2/(2(1-r^2)\varepsilon)]}
{2|\log r|}-1\right\rceil\right\}.
\]
La mesure optimale binaire agit sur le sous-espace engendré par les deux
branches de $B_R$. Elle peut être réalisée en annulant sur Bob les rotations
contrôlées par $B_0$, puis en mesurant $B_0$ dans la base $X$.

\section{Un parent à gap, et ses limites physiques}
Dans la représentation produit des qubits d'environnement, posons
\[
H_0=(I-\ket{\psi_p}\bra{\psi_p})_{AB_0}
+\sum_{j\geq1}\ket1\bra1_{E_j},
\quad \ket{\psi_p}=\sqrt p\ket{00}+\sqrt q\ket{11}.
\]
C'est l'opérateur positif de nombre d'excitations augmenté du projecteur
sur l'orthogonal de $\psi_p$ ; son état fondamental est unique et son gap
vaut un. Définir $H=UH_0U^*$ sur $U\mathcal D(H_0)$ conserve exactement
ce spectre et donne l'état précédent comme état fondamental.
Chaque différence $U\ket1\bra1_{E_j}U^*-\ket1\bra1_{E_j}$ est supportée
sur $\{A,E_j\}$ et sa norme est au plus $2\theta_j$. Le terme sur $AB_0$
devient quasi-local, avec erreur de troncature au plus
$2\sum_{j>R}\theta_j$. La perturbation totale est donc bornée, en norme,
par $2+2\sum_j\theta_j$ ; la description du parent est bien définie.

Ce modèle à interactions sommables est inhomogène, sans portée finie
et sans ordre topologique établi. Il est relié par un unitaire quasi-local
à un état produit par blocs. Surtout, la partition $A,B_R,C_R$ avec un
qubit Alice fixe n'est pas la géométrie de cônes de la proposition 51.2
de la note initiale. Les résultats ci-dessus \emph{ne démontrent pas}
l'optimalité de sa borne $2f(t)$, ni une réciproque de dualité de Haag.

\section{Antériorité, contrôles et clôture}
Les états à branches et les produits de facteurs de décohérence sont
classiques dans les travaux de Blume-Kohout et Zurek\ \cite{darwinism}.
La construction des POVM à marginale fixée relève de HJW\ \cite{hjw}.
La géométrie des états accessibles et les contraintes d'ensembles sont
étudiées par Jevtic et al.\ \cite{ellipsoid}. Les liens structurels entre
steering, relèvement et espérances conditionnelles\ \cite{lifting}, ainsi
que les liens avec les secteurs anyoniques\ \cite{anyons}, précèdent
cette note. Ogata\ \cite{ogata} fournit la notion de dualité relative
approchée utilisée dans la première partie de la piste.

Une recherche ciblée le 17 novembre 2026, détaillée dans
\texttt{piste51-cloture.md}, n'a pas identifié l'égalité minimax
\eqref{spatialexact} sous cette définition exacte. Cela ne certifie aucune
priorité : c'est une conséquence courte de structures connues.
Le script \texttt{verify51.py} contrôle indépendamment les préparations par
trace partielle, les témoins, des cas rationnels, les queues spatiales
et le spectre de petits parents. Il ne remplace pas les preuves et ne
certifie pas la nouveauté.

\textbf{Verdict.} L'approfondissement demandé est terminé : l'obstruction
mixte de la note initiale, la borne géométrique pure restreinte et le modèle
exact ci-dessus donnent un dossier vérifiable. Il reste une note technique
potentiellement utile, mais aucune découverte substantielle dont l'originalité
soit établie. La question générale de dualité quantitative reste ouverte.

\begin{thebibliography}{9}
\bibitem{hjw} L. P. Hughston, R. Jozsa, W. K. Wootters,
A complete classification of quantum ensembles having a given density matrix,
Phys. Lett. A 183 (1993), 14--18.
\href{https://doi.org/10.1016/0375-9601(93)90880-9}{doi:10.1016/0375-9601(93)90880-9}.
\bibitem{ellipsoid} S. Jevtic, M. F. Pusey, D. Jennings, T. Rudolph,
Quantum Steering Ellipsoids, Phys. Rev. Lett. 113, 020402 (2014).
\href{https://arxiv.org/abs/1303.4724}{arXiv:1303.4724}.
\bibitem{darwinism} R. Blume-Kohout, W. H. Zurek,
Quantum Darwinism: Entanglement, branches, and the emergent classicality
of redundantly stored quantum information, Phys. Rev. A 73, 062310 (2006).
\href{https://arxiv.org/abs/quant-ph/0505031}{arXiv:quant-ph/0505031}.
\bibitem{lifting} L. van Luijk et al.,
\href{https://arxiv.org/abs/2608.10783}{arXiv:2608.10783} (2026).
\bibitem{anyons} D. Wallick, H. Wilming,
\href{https://arxiv.org/abs/2609.01267}{arXiv:2609.01267} (2026).
\bibitem{ogata} Y. Ogata,
\href{https://arxiv.org/abs/2501.02398}{arXiv:2501.02398} (2025).
\end{thebibliography}
\end{document}
